\documentclass[12pt,a4paper]{article}

% ---------- Pacotes ----------
\usepackage[utf8]{inputenc}
\usepackage[T1]{fontenc}
\usepackage[a4paper,top=3cm,bottom=2cm,left=3cm,right=2cm]{geometry}
\usepackage{setspace}
\usepackage{indentfirst}
\usepackage{fancyhdr}
\usepackage{url}
\usepackage{hyperref}

\hypersetup{
    hidelinks,
    pdftitle={Resenha - O lugar das Humanidades na ideia de Universidade crítica},
    pdfauthor={Jeann Victor Batista},
    pdfsubject={Filosofia e Metodologia da Ciência}
}

\setstretch{1.5}
\setlength{\parindent}{1.25cm}
\setlength{\parskip}{0pt}
\setlength{\emergencystretch}{3em}

\pagestyle{fancy}
\fancyhf{}
\fancyfoot[C]{\thepage}
\renewcommand{\headrulewidth}{0pt}

% ---------- Dados do discente ----------
\newcommand{\NomeAluno}{Jeann Victor Batista}
\newcommand{\InstituicaoCurso}{Universidade Federal de Alfenas | Ciência da Computação}
\newcommand{\DisciplinaProfessor}{Filosofia e Metodologia da Ciência - José Francisco Lopes Xarão}

\begin{document}

% ======================================================
% CAPA / IDENTIFICAÇÃO
% ======================================================
\begin{center}
    \textbf{\NomeAluno}\footnote{Em conformidade com a Resolução nº 02/2025 (Diretrizes para Uso de Inteligência Artificial na UNIFAL-MG), declara-se que houve uso das ferramentas Claude, da Anthropic, como apoio à organização, à adequação ao modelo e à revisão linguística do texto. A leitura da obra, a seleção das ideias, a análise crítica e a responsabilidade pela versão entregue permanecem com o discente.}\\[2pt]
    \InstituicaoCurso\\[2pt]
    \DisciplinaProfessor\\[24pt]

    \textbf{\Large RESENHA}\\[6pt]
    \textit{O lugar das Humanidades na ideia de Universidade crítica}\\[4pt]
    \textit{\small The place of the Humanities in the idea of the critical University}
\end{center}

\vspace{18pt}

% ======================================================
% REFERÊNCIA DA OBRA
% ======================================================
\noindent\textbf{Resenha da obra:}\par
\noindent\setlength{\parindent}{0pt}%
FRAGA, Paulo Denisar. Lugar das Humanidades na ideia de Universidade crítica. In: FERRAREZI JR., Celso (Org.). \textit{A identidade docente no Ensino Superior e a universidade brasileira: uma contribuição da Universidade Federal de Alfenas ao debate nacional}. São Paulo; Alfenas: Scortecci; Unifal-MG, 2011. p. 177--195.
\setlength{\parindent}{1.25cm}

\vspace{12pt}

% ======================================================
% RESUMO TÉCNICO-CIENTÍFICO
% ======================================================
\noindent\textbf{Resumo}\par
\noindent\setlength{\parindent}{0pt}%
Esta resenha analisa o texto ``Lugar das Humanidades na ideia de Universidade crítica'', de Paulo Denisar Fraga, publicado como capítulo do livro \textit{A identidade docente no Ensino Superior e a universidade brasileira}. O tema do capítulo é a importância das Humanidades para que a Universidade participe da formação de sujeitos críticos. O problema discutido pelo autor consiste em saber se essa formação depende apenas de métodos de ensino e da atuação individual dos professores ou se exige uma concepção de Universidade que integre Ciências Empíricas e reflexão humanística. Fraga defende a segunda posição e recupera o modelo de Wilhelm von Humboldt, no qual ensino, pesquisa e formação geral não devem ser separados. Para desenvolver o argumento, o autor adota uma exposição ensaística, recorrendo a exemplos como a metáfora da casa, a relação entre ciência e responsabilidade no caso de Einstein e o surgimento da Bioética. A apreciação crítica relaciona a metáfora da casa ao romance \textit{1984}, de George Orwell, e aproxima a discussão da formação em Ciência da Computação. Considera-se que o capítulo apresenta uma defesa relevante das Humanidades, embora pudesse mostrar com maior precisão como a integração proposta seria realizada no cotidiano dos cursos universitários.

\vspace{6pt}
\noindent\textbf{Palavras-chave:} Universidade. Humanidades. Formação crítica. Ciência. Humboldt.
\setlength{\parindent}{1.25cm}

\vspace{12pt}

% ======================================================
% ABSTRACT
% ======================================================
\noindent\textbf{Abstract}\par
\noindent\setlength{\parindent}{0pt}%
\textit{This review analyzes the text ``The place of the Humanities in the idea of the critical University'', by Paulo Denisar Fraga, published as a chapter of the book A identidade docente no Ensino Superior e a universidade brasileira. The chapter addresses the importance of the Humanities in enabling the University to contribute to the formation of critical subjects. The author discusses whether such formation depends only on teaching methods and individual professors or requires a conception of the University that integrates empirical sciences and humanistic reflection. Fraga supports the second position and revisits Wilhelm von Humboldt's model, in which teaching, research, and general education should not be separated. His essayistic argument uses examples such as the metaphor of the house, the relationship between science and responsibility in Einstein's case, and the emergence of Bioethics. The critical assessment relates the metaphor of the house to George Orwell's novel 1984 and connects the discussion to education in Computer Science. The chapter offers a relevant defense of the Humanities, although it could explain more precisely how the proposed integration might be implemented in everyday university courses.}

\vspace{6pt}
\noindent\textit{\textbf{Keywords:} University. Humanities. Critical education. Science. Humboldt.}
\setlength{\parindent}{1.25cm}

\vspace{18pt}

% ======================================================
% RESENHA - CORPO CONTÍNUO, SEM SUBSEÇÕES
% ======================================================
\begin{center}
\textbf{\large Resenha}
\end{center}
\vspace{6pt}

Esta é uma resenha do texto ``Lugar das Humanidades na ideia de Universidade crítica'', escrito por Paulo Denisar Vasconcelos Fraga e publicado no livro \textit{A identidade docente no Ensino Superior e a universidade brasileira: uma contribuição da Universidade Federal de Alfenas ao debate nacional}, organizado por Celso Ferrarezi Jr. O capítulo teve origem em uma apresentação realizada na mesa-redonda ``A responsabilidade da universidade na formação do sujeito crítico'', durante o VII Seminário sobre Leitura e Produção no Ensino Superior, ocorrido na Unifal-MG em setembro de 2010. Essa origem ajuda a compreender o tom do texto, que mantém a forma de uma reflexão apresentada ao público, sem seguir uma exposição rígida ou uma pesquisa empírica.

Paulo Denisar Fraga é licenciado em Filosofia pela Universidade Federal de Santa Maria e mestre em Filosofia pela Universidade Estadual de Campinas. Foi o primeiro diretor do Instituto de Ciências Humanas e Letras da Universidade Federal de Alfenas, instituição em que atua como professor. Sua trajetória está ligada à Teoria Crítica, à tradição dialética e à discussão sobre a ideia de Universidade. Essas referências aparecem no capítulo por meio de autores como Marx, Adorno, Benjamin e Habermas, mas o texto não se limita a expor teorias: seu objetivo é discutir o lugar que as Humanidades ocupam na formação universitária (UNIFAL-MG, 2013).

O tema principal é a relação entre as Humanidades e a formação do sujeito crítico. O problema apresentado por Fraga pode ser resumido em uma pergunta: essa formação depende apenas de bons métodos didáticos e do esforço de professores considerados excelentes ou está ligada à própria concepção de Universidade? O autor rejeita a primeira possibilidade. Sua hipótese é que a responsabilidade de formar sujeitos críticos depende, antes de tudo, do modelo de Universidade adotado e do espaço reservado à reflexão humanística dentro dele (FRAGA, 2011, p. 177--179).

O capítulo possui quatro partes. Na primeira, Fraga apresenta a hipótese e informa que não seguirá uma exposição inteiramente sistemática. Em vez disso, reúne imagens e passagens que permitam observar o problema por diferentes ângulos. Essa escolha combina com o próprio conteúdo defendido, pois a crítica exige que o leitor não permaneça preso a apenas uma forma de compreender a realidade. O texto assume, assim, um caráter ensaístico: os exemplos não funcionam como etapas de uma demonstração fechada, mas como aproximações que reforçam a mesma tese.

Na segunda parte, o autor discute a relação entre Ciência e reflexão humanística. Um dos pontos iniciais é a diferença de prestígio entre essas áreas. A Ciência costuma aparecer como moderna, produtiva e útil, enquanto as Humanidades são frequentemente tratadas como antigas ou pouco práticas. Fraga questiona essa separação. Para ele, o conhecimento técnico é capaz de explicar como muitos processos acontecem, mas não responde sozinho por que determinada ação deve ser realizada, quais interesses ela atende ou quais consequências sociais pode produzir.

A metáfora da casa, retomada a partir de Bachelard e Fourez, é um dos exemplos mais claros do capítulo. Viver apenas no ``apartamento'' significa permanecer em um único plano, cercado por respostas prontas e por uma linguagem restrita à execução de tarefas. Descer ao porão ou subir ao sótão representa a possibilidade de mudar de perspectiva, recuperar experiências, formular novas perguntas e buscar razões que não aparecem no funcionamento imediato das coisas. As Humanidades cumprem essa função ao interromper a aceitação automática do conhecimento e perguntar por seus fundamentos, limites e finalidades (FRAGA, 2011, p. 180--182).

Fraga também recorre à diferença entre o atomismo de Demócrito e o de Epicuro, retomada pelo jovem Marx. A alteração espontânea no movimento dos átomos proposta por Epicuro permite ao autor discutir liberdade e subjetividade, mostrando que até uma questão aparentemente restrita à explicação da natureza pode envolver diferentes concepções de ser humano e de mundo. O exemplo reforça que as fronteiras entre Ciências Empíricas e Humanidades não são tão fechadas quanto parecem.

O caso de Einstein torna essa discussão mais concreta. O conhecimento da Física permitiu compreender processos necessários ao desenvolvimento da bomba atômica, mas não determinou como esse conhecimento deveria ser utilizado. O cientista participa da produção de uma descoberta sem controlar todas as decisões políticas e militares tomadas a partir dela. A Bioética aparece pelo mesmo motivo: o avanço das Ciências Biológicas ampliou aquilo que pode ser feito, mas também criou situações em que saber como realizar um procedimento não basta para decidir se ele deve ser realizado. Nesses casos, a reflexão ética não substitui a ciência; ela questiona os objetivos e as consequências de sua aplicação (FRAGA, 2011, p. 183--187).

Na terceira parte, o autor apresenta a relação entre Ciências e Humanidades no projeto moderno de Universidade. A principal referência é Wilhelm von Humboldt e a organização da Universidade de Berlim. Nesse modelo, ensino e pesquisa deveriam permanecer unidos. A Ciência não seria um conteúdo terminado, transmitido por quem sabe a quem apenas recebe, mas uma busca que envolve professores e estudantes. A Filosofia teria a tarefa de relacionar saberes especializados e impedir que a fragmentação das disciplinas eliminasse a reflexão sobre o conjunto (FRAGA, 2011, p. 188--191).

Esse modelo se opõe à ideia de Universidade reduzida à preparação imediata para o mercado. Quando o ensino superior é avaliado apenas pela rapidez com que fornece mão de obra ou resultados, atividades como pesquisa, extensão e formação cultural parecem secundárias. A consequência é que o estudante pode aprender técnicas eficientes sem compreender o contexto em que serão usadas. A defesa das Humanidades feita por Fraga não exige que todos abandonem a especialização. O que ela exige é que a especialização não seja confundida com uma formação completa.

Na parte final, Fraga reconhece que as Humanidades também possuem contradições e não ocupam automaticamente uma posição crítica. Elas podem reproduzir o mesmo pensamento positivista e fragmentado que pretendem questionar. Por isso, o autor não defende as Humanidades apenas por causa das disciplinas que recebem esse nome, mas pela função de relacionar conhecimentos, recuperar a dimensão histórica e imaginar possibilidades diferentes do presente. Ao retomar Walter Benjamin e parafrasear Sérgio Paulo Rouanet, conclui que o ``não-lugar'' da Filosofia, da História e da Literatura significa também a perda do pensamento questionador, da consciência histórica e da imaginação (FRAGA, 2011, p. 191--194).

O aspecto que mais chamou minha atenção foi a metáfora do apartamento, do porão e do sótão. Ela pode ser aproximada do romance \textit{1984}, de George Orwell, no qual o controle dos indivíduos depende também da limitação da linguagem, da memória e da possibilidade de questionar aquilo que é apresentado como verdade. As duas obras não tratam do mesmo problema, mas ajudam a perceber como uma visão única da realidade reduz a autonomia. Na metáfora discutida por Fraga, sair do apartamento significa reconhecer que há outros pontos de vista; no romance de Orwell, impedir essa saída é uma das formas de manter o poder (ORWELL, 2009).

Essa reflexão também se aplica à formação em Ciência da Computação. Boa parte do curso é dedicada a aprender como resolver problemas, otimizar algoritmos e construir sistemas. Esse conhecimento é necessário, mas não responde sozinho para quem uma tecnologia é desenvolvida, quais dados ela utiliza ou quais grupos podem ser prejudicados por suas decisões. Algoritmos de recomendação e sistemas de inteligência artificial, por exemplo, podem alcançar grande eficiência técnica e ainda reproduzir critérios que seus usuários não conhecem. A questão não é diminuir o valor da técnica, mas perceber que o ``como'' de um sistema precisa ser acompanhado por perguntas sobre seu propósito e suas consequências.

Nesse ponto está uma contribuição importante do capítulo. Fraga mostra que as Humanidades não são um complemento decorativo colocado ao lado das Ciências. Elas ajudam a avaliar os projetos que orientam a produção do conhecimento. Para estudantes de áreas técnicas, isso significa compreender que uma solução possível não é necessariamente uma solução desejável. A formação crítica aparece quando o profissional consegue analisar o problema que recebeu, em vez de apenas executar a tarefa da maneira mais eficiente.

Apesar da relevância do argumento, o texto apresenta uma limitação. Fraga defende de forma convincente a integração entre Humanidades e Ciências Empíricas, mas oferece poucas indicações de como ela poderia ocorrer na organização concreta dos cursos. Os exemplos de Einstein e da Bioética demonstram por que essa aproximação é necessária, porém não mostram como superar currículos fragmentados, como dividir responsabilidades entre departamentos ou como avaliar uma formação realmente interdisciplinar. A análise ganharia força com experiências universitárias específicas, principalmente da realidade brasileira, que mostrassem os resultados e as dificuldades dessa integração.

Essa observação não anula o argumento central. Ao contrário, indica um caminho para continuar a discussão iniciada pelo autor. Se a concepção de Universidade é mais importante do que iniciativas isoladas, não basta incluir uma disciplina de Filosofia em um currículo técnico e considerar o problema resolvido. Seria necessário criar oportunidades em que questões éticas, históricas e sociais participassem da própria análise de problemas científicos e tecnológicos.

A leitura é recomendada a estudantes de graduação, professores e gestores universitários, especialmente àqueles que discutem currículos e a função social do ensino superior. Também é pertinente para estudantes das Engenharias e da Ciência da Computação, porque questiona uma formação limitada ao domínio de ferramentas. O principal mérito do capítulo está em mostrar que a Universidade forma sujeitos críticos quando trata o conhecimento como busca e reflexão, e não apenas como informação pronta ou treinamento profissional. O texto cumpre parte dessa proposta ao levar o leitor a examinar a própria experiência universitária e perguntar que tipo de formação está recebendo.

% ======================================================
% REFERÊNCIAS
% ======================================================
\begin{center}
\textbf{\large Referências}
\end{center}
\vspace{6pt}

\begin{singlespace}
\sloppy
\noindent\setlength{\parindent}{0pt}%
FRAGA, Paulo Denisar. Lugar das Humanidades na ideia de Universidade crítica. In: FERRAREZI JR., Celso (Org.). \textit{A identidade docente no Ensino Superior e a universidade brasileira: uma contribuição da Universidade Federal de Alfenas ao debate nacional}. São Paulo; Alfenas: Scortecci; Unifal-MG, 2011. p. 177--195.

\par\vspace{6pt}\noindent
ORWELL, George. \textit{1984}. São Paulo: Companhia das Letras, 2009.

\par\vspace{6pt}\noindent
UNIVERSIDADE FEDERAL DE ALFENAS. Diretoria do Instituto de Ciências Humanas e Letras da UNIFAL-MG é reeleita. \textit{Assessoria de Comunicação Social}, Alfenas, 12 abr. 2013. Disponível em: \url{https://www.unifal-mg.edu.br/comunicacao/diretoriaichlreeleita}. Acesso em: 3 out. 2026.
\end{singlespace}

\end{document}

